\section{Discussion and Next Tests}
\label{sec:discussion}
The evidence supports a specific conclusion: with a common fitted forecaster and pilot-only inputs, the current OEH feature set improves prediction of the registered learner-specific censored requirement on new instances of the studied stochastic MDP families. The seventh cohort passes all 22 registered forecast contrasts with positive simultaneous bounds. It does not establish that every new component is independently necessary. Development-stage feature arms informed the frozen candidate, but a clean, same-forecaster component ablation on a new independent benchmark remains to be completed. Any new ablation on the already opened seventh cohort is exploratory.

The privileged TQ/SACo augmentation shows why information access matters. It uses author-code measures and reference datasets produced from the true MDP, whereas operational OEH sees only a pilot. The fourth cohort's tier-2 failure and the seventh cohort's positive tier-2 result both remain visible. The selection experiment compares with seven fixed learner implementations, including author BCQ and SPIBB code and a pinned CQL library implementation; it does not compare against every published policy-selection system on the same task. The registered internal v5 comparator has slightly higher mean return and no demonstrated disadvantage. Accordingly, neither endpoint supports a literature-wide state-of-the-art claim. Better scores also do not establish cheaper deployment until posterior-simulation time, diagnostic cost, and downstream data savings are measured.

Support restrictions create a fundamental identification limit. If two MDPs agree on every logged action but assign different rewards to an unseen action, a pilot-only statistic has the same input in both and cannot identify the optimal action in both. The posterior convention for unobserved rewards is therefore a modeling choice, not a finite-sample safety certificate. The weak data-only OEH-H rank association and raw OEH-B miscalibration reinforce this limitation.

The published FourRooms component test has a NO-GO verdict against maximum likelihood and uses partial visual observations, so it does not close the full-method transfer gap. The next confirmatory study should freeze a minimal component ablation and transfer protocol before running an independently implemented stochastic tabular benchmark. It should use the same learner, budget, behavior, and endpoint definitions where possible; report paired MDP-level intervals and all unsuccessful cases; and measure diagnostic time and downstream data cost separately. If transfer fails, the present within-family forecasting result remains the supported scope. The paper and thesis retain prior negative gates regardless of the next test.
